% ===== Block G: PDF p.31–34 | 原文 §3.2（Givental's action） | 唯一编号公式 3.25–3.48 =====
% [NOTE: 本块自 p.31 页内的第二条项目符号（(3.25) 所在的“固定目标 Kähler 时空”一条）开始；该页顶部 Example 3.2 的收尾句已由 Block F 译完并闭合环境。原文的 §3.2 标题排在 p.31 中部（在本条项目符号之后），按 Lead 裁决由本块就地写入 \subsection；原文 §3.2 内部的三个小标题 3.2.1/3.2.2/3.2.3 亦由本块就地以 \subsubsection 写入。]
% [NOTE: (3.26)、(3.35)、(3.39) 在原书中是稳定图的图示而非代数式；按本译注本“不画图”的统一规定（与 Block F 处理 (3.9)-(3.12) 的办法一致），改用文字描述图的结构与标号，编号照原文保留。]
% [原文 p.31]

\begin{itemize}
\item \emph{拓扑弦振幅}（Topological string amplitudes）。固定目标 Kähler 时空 $(X,\omega)$ 上的拓扑弦振幅，恰好是某个 CohFT 的关联函数；该 CohFT 底层的相空间是分次向量空间
\begin{equation}
  V=\bigoplus_{\beta\in H_{2}(X,\Z)}H^{\bullet}(X,\Z)\cdot q^{\beta}\,,\qquad
  \eta(\gamma_{1},\gamma_{2})=\int_{X}\gamma_{1}\smile\gamma_{2}\,.
  \tag{3.25}
\end{equation}
% [Tutor Note: (3.25) 的指标写的是同调 $H_{2}(X,\Z)$：曲线类 $\beta$ 取在目标空间的二阶同调里（Poincaré 对偶下等同于 $H^{2}(X,\Z)$ 中的类），而每个分次片里的系数取上同调 $H^{\bullet}(X,\Z)$。核对：$\beta=0$ 的分次片是 $H^{\bullet}(X,\Z)\cdot q^{0}=H^{\bullet}(X,\Z)$，包含单位 $\mathbf{1}\in H^{0}(X,\Z)$；$\eta$ 是目标上的上同调杯积配对，故 $\eta$ 非退化且对称。这也正是原文紧接着说的“$V$ 无穷维、但每个分次片有限维”。]

这里 $q$ 是形式变量，称为 Novikov 变量，定义为 $q^{\beta}=\ee^{-\int_{\beta}\omega}$，用来按曲线类 $\beta\in H_{2}(X,\Z)$ 的辛面积对其贡献分次。在此语境下，相空间 $V$ 是无穷维的，但由 $H_{2}(X,\Z)$ 分次，各分量 $H^{\bullet}(X,\Z)$ 是有限维的。上同调中的单位 $\mathbf{1}\in H^{0}(X,\Z)$ 充当相应 CohFT 的单位。这一构造是 Kontsevich 与 Manin 对 CohFT 作公理化定义时的动机性例子，见 [KM94]。
% [Tutor Note: Novikov 变量 $q^{\beta}=\ee^{-\int_{\beta}\omega}$ 只是形式记号（$\beta\mapsto\int_{\beta}\omega$ 未必取离散值），关键在于它把辛面积放进指数，从而 $q^{\beta}q^{\beta'}=q^{\beta+\beta'}$：核对方法是对 $\int_{\beta+\beta'}\omega=\int_{\beta}\omega+\int_{\beta'}\omega$ 取指数。指数上的负号纯属约定，换成 $q^{\beta}=\ee^{+\int_{\beta}\omega}$ 只会把分次的符号反过来。]
\end{itemize}

\subsection{Givental 作用（Givental's action）}

我们已经看到 $\Mc_{g,n}$ 如何呈现递归的边界结构。一个自然的问题是：能否利用这样的递归结构来定义或计算 CohFT？答案是肯定的，其根源在于 A. Givental 关于拓扑弦论中局部化计算 [GP99] 的著作 [Giv01a; Giv01b]。更确切地说，Givental 在 CohFT 上定义了两个作用：旋转作用与平移作用。

\subsubsection{Rotation 旋转}

对固定的 $(g,n)$，我们有参数化 $\Mc_{g,n}$ 的边界的全部可能稳定图的清单。若给定 $(V,\eta)$ 上的一个 CohFT $\Omega$，自然的做法是用 $\Omega$ 提供的上同调类装饰所有顶点，从而得到 $\Mc_{\Gamma}$ 上的一个上同调类。例如：
\begin{equation}
  \parbox{0.62\textwidth}{\centering\small\kaishu 〔原文此处为图示：一个亏格为 $2$ 的顶点（圆圈内标 $2$），其右侧带一个自环，自环的两条半边分别标 $\alpha$（上）、$\beta$（下）；左侧伸出两条leg，分别标 $\nu_{1}$（上）、$\nu_{2}$（下）。〕}
  \;\rightsquigarrow\;\Omega_{2;\nu_{1},\nu_{2},\alpha,\beta}\,.
  \tag{3.26}
\end{equation}
% [Tutor Note: (3.26) 左端这个图只有一个顶点，故顶点装饰就是 $\Omega_{2;\nu_{1},\nu_{2},\alpha,\beta}$，其下标按“leg、半边”的次序排列；自环贡献因子 $E^{\alpha,\beta}$，两条leg分别贡献 $L^{\nu_{1}}_{\mu_{1}}$、$L^{\nu_{2}}_{\mu_{2}}$，合起来正是 (3.27)。核对缩并：在 (3.27) 中数指标，$\nu_{1},\nu_{2}$ 各出现两次（一次在 $\Omega$，一次在 $L$），$\alpha,\beta$ 各出现两次（一次在 $\Omega$，一次在 $E$），恰好吃掉全部内部指标，只剩叶子上的 $\mu_{1},\mu_{2}$。]

为在 $\Mc_{g,n}$ 上得到上同调类，我们必须用取值为上同调类的矩阵 $E^{\nu_{h},\nu_{h'}}$（先验地任意）缩并边上所有的指标，用取值为上同调类的矩阵 $L^{\nu_{i}}_{\mu_{i}}$（先验地任意）缩并叶子处的指标，并把结果经粘合映射 $\xi_{\Gamma}$ 前推。在上述例子中，我们会得到
\begin{equation}
  \Omega_{2;\nu_{1},\nu_{2},\alpha,\beta}E^{\alpha,\beta}L^{\nu_{1}}_{\mu_{1}}L^{\nu_{2}}_{\mu_{2}}\,.
  \tag{3.27}
\end{equation}
其中 $\mu_{i}$ 表示第 $i$ 片叶子（即第 $i$ 个标记点）处一个固定的装饰。

除以自然的自同构因子并对所有可能的稳定图求和，我们得到如下形式的表达式
\begin{equation}
  \sum_{\Gamma\ \mathrm{of\ type}\ (g,n)}\frac{1}{|\Aut(\Gamma)|}\,\xi_{\Gamma,*}\left(\prod_{v\in V(\Gamma)}\Omega_{g(v);(\nu_{h})_{h\rightsquigarrow v}}\right)\left(\prod_{e=(h,h')\in E(\Gamma)}E^{\nu_{h},\nu_{h'}}\right)\left(\prod_{i=1}^{n}L^{\nu_{i}}_{\mu_{i}}\right)\,.
  \tag{3.28}
\end{equation}
这里 $h\rightsquigarrow v$ 表示与顶点 $v$ 相关联的任意半边 $h$。
% [Tutor Note: 因子 $1/|\Aut(\Gamma)|$ 是必要的：粘合映射 $\xi_{\Gamma}$ 把 $\prod_{v\in V(\Gamma)}\Mc_{g(v),n(v)}$ 打到 $\Mc_{g,n}$，其一般纤维的大小恰为 $|\Aut(\Gamma)|$，故要按自同构因子作加权（有理系数）。核对：对 (3.26) 的图，自环的两条半边可以互换而图不变（叶子 $\nu_{1},\nu_{2}$ 的标号不同，不能互换），故 $|\Aut(\Gamma)|=2$；这也说明 (3.27) 中不出现 $1/2$ 是因为那里尚未作自同构加权。]
% [Tutor Note: 记号切换：本页末尾起，原文把 (3.28) 中的 $E^{\nu_{h},\nu_{h'}}$、$L^{\nu_{i}}_{\mu_{i}}$ 改写成一般的矩阵 $E^{\mu,\nu}$、$L^{\nu}_{\mu}$，用以强调在整顿记号之前它们是“先验地任意”的上同调类值矩阵；到 (3.29)-(3.30) 才把它们具体化为关于 $u,v$ 的幂级数。读式子时只需注意指标字母的位置（上/下）而非字母本身。]

自然的问题是：由 (3.28) 得到的这族上同调类何时构成一个 CohFT？结果表明 (3.28) 过于朴素：矩阵 $E^{\mu,\nu}$ 与 $L^{\nu}_{\mu}$ 不能是任意的，而应涉及 $\psi$ 类的特定组合。这一条件由单个元素刻画，称为旋转矩阵。
% [Tutor Note: 原文此句作 “specific combinations $\psi$-classes”，脱漏了介词 of，应为 “specific combinations of $\psi$-classes”；译文按“$\psi$ 类的特定组合”处理，笔误已登记入 ERRATA。]
% [原文 p.32]

$(V,\eta)$ 上的一个\emph{旋转矩阵}是一个 $\mathrm{End}(V)$ 值的幂级数，它在次数 $0$ 处为恒等映射，并且关于 $\eta$ 满足辛条件：
\begin{equation}
  R^{\nu}_{\mu}(u)=\delta^{\nu}_{\mu}+\sum_{k\ge1}(R_{k})^{\nu}_{\mu}u^{k}\in\Q[[u]]\,,\qquad
  R^{\mu}_{\alpha}(u)\,\eta^{\alpha,\beta}\,R^{\nu}_{\beta}(-u)=\eta^{\mu,\nu}\,.
  \tag{3.29}
\end{equation}
% [Tutor Note: 第二个等式缩并的是逆配对 $\eta^{\alpha,\beta}$（上标），其含义是 $R(u)$ 保持 $\eta^{-1}$；等价写法是 $R(u)^{-1}=\eta^{-1}R(-u)^{T}\eta$，这就是“辛条件”一词的来历。核对方法：把 $R^{\nu}_{\mu}(u)=\delta^{\nu}_{\mu}+(R_{1})^{\nu}_{\mu}u+O(u^{2})$ 代入并比较 $u$ 的一次项，得 $(R_{1})^{\mu}{}_{\alpha}\eta^{\alpha,\nu}=\eta^{\mu,\beta}(R_{1})^{\nu}{}_{\beta}$，即 $R_{1}\eta^{-1}$ 是对称矩阵（等价地，$R_{1}$ 关于 $\eta^{-1}$ 自伴随）；取 $R(u)=\mathrm{id}$（即 $R_{1}=0$）时等式退化为 $\eta^{\mu,\nu}=\eta^{\mu,\nu}$。]

给定旋转矩阵，定义边装饰为如下两变量的 $V^{\otimes 2}$ 值幂级数\footnote{注意，若干作者用 $R^{-1}$ 代替 $R$。这里我们遵循 Givental 的约定。}：
\begin{equation}
  E^{\mu,\nu}(u,v)=\frac{\eta^{\mu,\nu}-R^{\mu}_{\alpha}(u)\,\eta^{\alpha,\beta}\,R^{\nu}_{\beta}(v)}{u+v}\in\Q[[u,v]]\,.
  \tag{3.30}
\end{equation}
% [Tutor Note: 正则性就是“分子被 $u+v$ 整除”：记分子为 $f(u,v)$，由 (3.29) 取 $v=-u$ 立得 $f(u,-u)=\eta^{\mu,\nu}-R^{\mu}_{\alpha}(u)\eta^{\alpha,\beta}R^{\nu}_{\beta}(-u)=0$，故 $f$ 含因子 $u+v$。顺带算出常数项（后面 (3.37)、(3.41)、(3.42) 要用）：把一阶展开代入可见 $f$ 的一次部分为 $-(R_{1})^{\mu}{}_{\alpha}\eta^{\alpha,\nu}(u+v)$，故 $E^{\mu,\nu}_{0,0}=-(R_{1})^{\mu}{}_{\alpha}\eta^{\alpha,\nu}$，一般并不等于 $\eta^{\mu,\nu}$。]

辛条件保证 $E^{\mu,\nu}(u,v)$ 沿 $u+v=0$ 是正则的。通过展开 $E^{\mu,\nu}(u,v)=\sum_{k,\ell\ge0}E^{\mu,\nu}_{k,\ell}u^{k}v^{\ell}$ 定义标量 $E^{\mu,\nu}_{k,\ell}$。

\begin{definition}[3.3]
考虑 $(V,\eta)$ 上的一个 CohFT $\Omega$ 以及一个旋转矩阵 $R$。我们定义一族新的、取值于上同调类的线性映射
\begin{equation}
  R\Omega_{g,n}:V^{\otimes n}\longrightarrow H^{2\bullet}(\Mc_{g,n},\Q)
  \tag{3.31}
\end{equation}
其具体构造如下。对每个 $(g,n)$ 型的稳定图 $\Gamma$，通过以下构造定义它的一个贡献：
\begin{itemize}
\item 在 $\Gamma$ 的每个顶点 $v$ 处放置 $\Omega_{g(v);(\nu_{h})_{h\rightsquigarrow v}}$，其中与 $v$ 相连的半边处带有任意的装饰 $\nu_{h}$；
\item 在 $\Gamma$ 的第 $i$ 片叶子处放置 $R^{\nu_{i}}_{\mu_{i}}(\psi_{i})$；
\item 在 $\Gamma$ 的每条边 $e=(h,h')$ 处放置 $E^{\nu_{h},\nu_{h'}}(\psi_{h},\psi_{h'})$；
\item 缩并所有指标。
\end{itemize}
% [Tutor Note: (3.31) 的目标写成 $H^{2\bullet}$ 而不是 $H^{\bullet}$：$R$ 与 $E$ 的系数矩阵不改变 cocycle 的偶次性，故旋转后的类仍在偶次上同调中。核对：与 Block F 的 (3.4)、(3.22) 以及本块的 (3.44) 对照即可；若某个 CohFT 本身带混合次数，此处的次数记账要相应调整。]
换言之，我们得到上同调类
\begin{equation}
  \mathrm{Contr}_{\Gamma;\mu_{1}\dots\mu_{n}}=\left(\prod_{v\in V(\Gamma)}\Omega_{g(v);(\nu_{h})_{h\rightsquigarrow v}}\right)\left(\prod_{e=(h,h')\in E(\Gamma)}E^{\nu_{h},\nu_{h'}}(\psi_{h},\psi_{h'})\right)\left(\prod_{i=1}^{n}R^{\nu_{i}}_{\mu_{i}}(\psi_{i})\right)\,.
  \tag{3.32}
\end{equation}
尽管表达式 $E^{\nu_{h},\nu_{h'}}(\psi_{h},\psi_{h'})$ 与 $R^{\nu_{i}}_{\mu_{i}}(\psi_{i})$ 先验地有无穷多项，但由于上同调次数上的理由，它们会截断。
% [Tutor Note: 截断的理由：$\psi_{h}$ 是次数 $2$ 的类，$E^{\nu_{h},\nu_{h'}}(\psi_{h},\psi_{h'})$ 的 $u^{k}v^{\ell}$ 项贡献次数 $2(k+\ell)$，把 $R$ 的各项的固定次数加上去后，超过 $\dim\Mc_{g,n}$ 的项在积分中为零。核对最简情形 $(g,n)=(0,3)$：$\dim\Mc_{0,3}=0$，故 $\psi_{i}=0$，(3.32) 只剩 $k=\ell=0$ 的常数项，这正是 (3.34) 的来源。]

定义 $R\Omega_{g;\mu_{1},\dots,\mu_{n}}$ 为所有稳定图的贡献之和，即在经前推到模空间并按自同构因子加权之后的和：
\begin{equation}
  R\Omega_{g;\mu_{1},\dots,\mu_{n}}=\sum_{\Gamma\ \mathrm{of\ type}\ (g,n)}\frac{1}{|\Aut(\Gamma)|}\,\xi_{\Gamma,*}\mathrm{Contr}_{\Gamma;\mu_{1},\dots,\mu_{n}}\,.
  \tag{3.33}
\end{equation}
% [Tutor Note: 记号对应：线性映射 $R\Omega_{g,n}$（(3.31)）作用在基上给出 $R\Omega_{g;\mu_{1},\dots,\mu_{n}}$，即 $R\Omega_{g,n}(e_{\mu_{1}}\otimes\cdots\otimes e_{\mu_{n}})=R\Omega_{g;\mu_{1},\dots,\mu_{n}}$，与 Block F 中 (3.4) 对 $\Omega$ 的约定一致。下文例题里两个记号混用（如 (3.34) 的 $R\Omega_{0,3}$ 指映射），读者按“带分号是分量、不带分号是映射”来区分即可。]
\end{definition}

让我们分析低拓扑中的一些例子。

% [原文 p.33]
\begin{itemize}
\item $R\Omega_{0,3}$。$(0,3)$ 型的稳定图只有一个，而且由于维数上的理由，叶子处的装饰 $R(\psi_{i})$ 就是恒等映射。于是我们得到
\begin{equation}
  R\Omega_{0,3}=\Omega_{0,3}\,.
  \tag{3.34}
\end{equation}
% [Tutor Note: 核对：$\dim\Mc_{0,3}=0$，故 $\psi_{1}=\psi_{2}=\psi_{3}=0$，$R(\psi_{i})=R(0)=\delta$ 正是恒等映射；此时唯一的稳定图没有边，其贡献就是顶点装饰 $\Omega_{0;\mu_{1},\mu_{2},\mu_{3}}=\Omega_{0,3}$。这也是“for dimensional reasons”的确切内容。]
\item $R\Omega_{0,4}$。$(0,4)$ 型的稳定图如下：
\begin{equation}
  \parbox{0.86\textwidth}{\centering\small\kaishu 〔原文此处为图示：左图为 $\Gamma_{0}$，即带四条leg的单个亏格 $0$ 顶点，四条leg自左上顺时针依次标 $1$、$4$、$3$、$2$（即 $1$ 左上、$4$ 右上、$2$ 左下、$3$ 右下）；右图为 $\Gamma_{ij|k\ell}$，即两个亏格 $0$ 顶点以一条边相连，左顶点伸出两条leg标 $i$（上）、$j$（下），右顶点伸出两条leg标 $k$（上）、$\ell$（下）。〕}
  \tag{3.35}
\end{equation}
其中 $ij|k\ell\in\{12|34,\,13|24,\,14|23\}$。稳定图 $\Gamma_{0}$ 的贡献由
\begin{equation}
\begin{aligned}
  \mathrm{Contr}_{\Gamma_{0};\mu_{1},\mu_{2},\mu_{3},\mu_{4}}
  &=\Omega_{0;\mu_{1},\mu_{2},\mu_{3},\mu_{4}}\\
  &\quad+\Omega_{0;\alpha,\mu_{2},\mu_{3},\mu_{4}}(R_{1})^{\alpha}_{\mu_{1}}\psi_{1}
  +\Omega_{0;\alpha,\mu_{1},\mu_{3},\mu_{4}}(R_{1})^{\alpha}_{\mu_{2}}\psi_{2}\\
  &\quad+\Omega_{0;\alpha,\mu_{1},\mu_{2},\mu_{4}}(R_{1})^{\alpha}_{\mu_{3}}\psi_{3}
  +\Omega_{0;\alpha,\mu_{1},\mu_{2},\mu_{3}}(R_{1})^{\alpha}_{\mu_{4}}\psi_{4}\,.
\end{aligned}
  \tag{3.36}
\end{equation}
给出。稳定图 $\Gamma_{ij|k\ell}$ 的贡献由
\begin{equation}
  \mathrm{Contr}_{\Gamma_{ij|k\ell};\mu_{1},\mu_{2},\mu_{3},\mu_{4}}
  =\Omega_{0;\mu_{i},\mu_{j},\alpha}\,E^{\alpha,\beta}_{0,0}\,\Omega_{0;\beta,\mu_{k},\mu_{\ell}}\,.
  \tag{3.37}
\end{equation}
给出。可以证明 $\xi_{\Gamma_{ij|k\ell},*}\mathbf{1}=[\Gamma_{ij|k\ell}]=\kappa_{1}$，于是我们得到
\begin{equation}
\begin{aligned}
  R\Omega_{0;\mu_{1},\mu_{2},\mu_{3},\mu_{4}}
  &=\Omega_{0;\mu_{1},\mu_{2},\mu_{3},\mu_{4}}\\
  &\quad+\Omega_{0;\alpha,\mu_{2},\mu_{3},\mu_{4}}(R_{1})^{\alpha}_{\mu_{1}}\psi_{1}
  +\Omega_{0;\alpha,\mu_{1},\mu_{3},\mu_{4}}(R_{1})^{\alpha}_{\mu_{2}}\psi_{2}\\
  &\quad+\Omega_{0;\alpha,\mu_{1},\mu_{2},\mu_{4}}(R_{1})^{\alpha}_{\mu_{3}}\psi_{3}
  +\Omega_{0;\alpha,\mu_{1},\mu_{2},\mu_{3}}(R_{1})^{\alpha}_{\mu_{4}}\psi_{4}\\
  &\quad+\left(\sum_{ij|k\ell}\Omega_{0;\mu_{i},\mu_{j},\alpha}E^{\alpha,\beta}_{0,0}\Omega_{0;\beta,\mu_{k},\mu_{\ell}}\right)\kappa_{1}\,.
\end{aligned}
  \tag{3.38}
\end{equation}
% [Tutor Note: (3.36) 中四个一次项都把被求和的指标 $\alpha$ 写在 $\Omega$ 下标的首位（例如 $\psi_{2}$ 项写作 $\Omega_{0;\alpha,\mu_{1},\mu_{3},\mu_{4}}$，而非按位置写作 $\Omega_{0;\mu_{1},\alpha,\mu_{3},\mu_{4}}$）。这不是笔误：CohFT 公理 (i)（对称性，见 Block F 中 (3.4) 之后的公理 (i)）保证 $\Omega_{0;\mu_{1},\mu_{2},\mu_{3},\mu_{4}}$ 对四个下标的任意置换不变，故两种写法给出同一个标量。核对：$\psi_{1}$ 项按位代入确实是 $\Omega_{0;\alpha,\mu_{2},\mu_{3},\mu_{4}}$，与原文一致。]
% [Tutor Note: (3.36) 与 (3.38) 的差别只在最后的 $\kappa_{1}$ 项：(3.38) 的前四行逐字等于 (3.36) 的右端（因为在 (3.34) 中 $\Gamma_{0}$ 的贡献就是 $\mathrm{Contr}_{\Gamma_{0}}$），而 $\sum_{ij|k\ell}$ 那项来自三个 $\Gamma_{ij|k\ell}$。核对：$\Mc_{0,4}\cong\PP^{1}$ 恰有三个边界点，对应 $\binom{4}{2}/2=3$ 个剖分 $12|34,13|24,14|23$；又 $|\Aut(\Gamma_{ij|k\ell})|=1$（两条leg分居两个顶点且标号不同），故 (3.37)-(3.38) 中没有额外的自同构因子。]
\item $R\Omega_{1,1}$。$(1,1)$ 型的稳定图有两个：
\begin{equation}
  \parbox{0.86\textwidth}{\centering\small\kaishu 〔原文此处为图示：左图 $\Gamma$ 是一条标 $1$ 的leg连到亏格 $1$ 的单个顶点（圈内标 $1$）；右图 $\Gamma'$ 是一条标 $1$ 的leg连到亏格 $0$ 的单个顶点（圈内标 $0$），该顶点还带一个自环。〕}
  \tag{3.39}
\end{equation}
$\Gamma$ 的贡献是
\begin{equation}
  \mathrm{Contr}_{\Gamma;\mu}=\Omega_{1;\mu}+\Omega_{1;\nu}(R_{1})^{\nu}_{\mu}\psi_{1}\,.
  \tag{3.40}
\end{equation}
对单环图 $\Gamma'$，我们得到
\begin{equation}
  \mathrm{Contr}_{\Gamma';\mu}=\Omega_{0;\mu,\alpha,\beta}E^{\alpha,\beta}_{0,0}\,.
  \tag{3.41}
\end{equation}
可以证明 $\frac{1}{2}\xi_{\Gamma',*}\mathbf{1}=[\Gamma']=12\psi_{1}$，于是
\begin{equation}
  R\Omega_{1;\mu}=\Omega_{1;\mu}+\left(\Omega_{1;\nu}(R_{1})^{\nu}_{\mu}+12\,\Omega_{0;\mu,\alpha,\beta}E^{\alpha,\beta}_{0,0}\right)\psi_{1}\,.
  \tag{3.42}
\end{equation}
% [Tutor Note: (3.42) 中 $12$ 与 (3.41) 前的 $\frac12$ 必须一起看：$|\Aut(\Gamma')|=2$（自环的两条半边可互换），故 $\frac{1}{|\Aut(\Gamma')|}\xi_{\Gamma',*}\mathbf{1}=\frac12[\Gamma']$；而 $[\Gamma']=12\psi_{1}$ 是 $\Mc_{1,1}$ 的边界类（$\Mc_{1,1}$ 只有一个边界点，其上 $\delta=12\lambda_{1}$）。数值核对：$\int_{\Mc_{1,1}}\psi_{1}=\frac{1}{24}$，故 $\int[\Gamma']=12\cdot\frac{1}{24}=\frac12$，正好等于该边界点在轨形意义下（结点曲线有 $|\Aut|=\Z/2$）所占的基本类比例。]
\end{itemize}

这一构造的要点在于：所得的那族取值于上同调类的映射 $R\Omega$ 构成一个 CohFT。

\begin{proposition}[3.4]
取值于上同调类的线性映射族 $R\Omega=(R\Omega_{g,n})_{2g-2+n>0}$ 构成 $(V,\eta)$ 上的一个 CohFT。此外，旋转作用构成一个右群作用。
\end{proposition}

% [原文 p.34]
\subsubsection{Translation 平移}

旋转作用通过一种 2 点关联函数（即旋转矩阵）把 CohFT 附着起来，从而利用了粘合映射。还有另一个我们可以纳入考虑的重言映射：遗忘映射。从图形上看，遗忘映射剪去图中的一片叶子，而这片叶子可以（在被遗忘之前）用一个 1 点关联函数来装饰。与旋转的情形一样，正确的做法是用 $\psi$ 类的特定组合来装饰被遗忘的叶子。平移作用正是把这一点纳入考虑。

一个\emph{平移}是一个在次数 $0$ 与 $1$ 处消失的 $V$ 值幂级数：
\begin{equation}
  T^{\mu}(u)=\sum_{d\ge1}(T_{d})^{\mu}u^{d+1}\in u^{2}\Q[[u]]\,.
  \tag{3.43}
\end{equation}
% [Tutor Note: 求和从 $d\ge1$ 起、且幂次写成 $u^{d+1}$，正是为了与“在次数 $0$ 与 $1$ 处消失”相容：$T$ 的最低次项是 $u^{2}$。核对：$d=1$ 项为 $(T_{1})^{\mu}u^{2}$，与“次数 $1$ 处消失”不矛盾（那里指 $u^{0}$ 与 $u^{1}$ 的系数为零）。这个错开一位与原文紧接着所说的求和有限性有关。]

\begin{definition}[3.5]
考虑 $(V,\eta)$ 上的一个 CohFT $\Omega$，以及一个平移 $T$。我们定义一族取值于上同调类的线性映射
\begin{equation}
  T\Omega_{g,n}:V^{\otimes n}\longrightarrow H^{2\bullet}(\Mc_{g,n},\Q)
  \tag{3.44}
\end{equation}
其方式是设定
\begin{equation}
  T\Omega_{g;\mu_{1},\dots,\mu_{n}}=\sum_{m\ge0}\frac{1}{m!}\pi_{m,*}\Omega_{g;\mu_{1},\dots,\mu_{n},\nu_{1},\dots,\nu_{m}}T^{\nu_{1}}(\psi_{n+1})\cdots T^{\nu_{m}}(\psi_{n+m})\,.
  \tag{3.45}
\end{equation}
\end{definition}

这里 $\pi_{m}:\Mc_{g,n+m}\to\Mc_{g,n}$ 是遗忘最后 $m$ 个标记点的映射。注意 $T$ 在次数 $0$ 与 $1$ 处消失，这保证了上述求和实际上是有限的。
% [Tutor Note: (3.45) 中 $m=0$ 的项就是 $\Omega_{g;\mu_{1},\dots,\mu_{n}}$ 本身（不平移）；因子 $\frac{1}{m!}$ 来自被遗忘的 $m$ 个标记点的置换对称性：$\pi_{m,*}$ 的纤维里这 $m$ 个点是互不区分的，按 $m!$ 个次序重复计数，故取平均。核对 $m=1$ 项：$\pi_{1,*}\Omega_{g;\mu_{1},\dots,\mu_{n},\nu}T^{\nu}(\psi_{n+1})$，恰好把多出的一个标记点遗忘掉。]

\begin{proposition}[3.6]
取值于上同调类的线性映射族 $T\Omega=(T\Omega_{g,n})_{2g-2+n>0}$ 构成 $(V,\eta)$ 上的一个 CohFT。此外，平移作用构成一个 Abel 群作用。
\end{proposition}

也可以验证旋转与平移之组合所满足的复合律。其结果与平面上旋转与平移的作用相平行，名称即由此而来。

\subsubsection{Examples and Teleman 例子与 Teleman 定理}

若干 CohFT 可由 Givental 作用表达。这里我们给出 Weil–Petersson 类与 Hodge 类的情形。

\begin{exercise}[3.2]
证明：$\exp(2\pi^{2}\kappa_{1})$ 是从平凡 CohFT 出发、经如下平移作用而得到的 CohFT：
\begin{equation}
  T(u)=-\sum_{k\ge1}\frac{(-2\pi^{2})^{k}}{k!}u^{k+1}=u\big(1-\ee^{-2\pi^{2}u}\big)\,.
  \tag{3.46}
\end{equation}
\end{exercise}
% [Tutor Note: (3.46) 中两个表达式的一致性可直接核对：$\sum_{k\ge1}\frac{(-2\pi^{2})^{k}}{k!}u^{k+1}=u(\ee^{-2\pi^{2}u}-1)$，再取负号即得 $u(1-\ee^{-2\pi^{2}u})$。逐项看，$k=1$ 给出左端 $2\pi^{2}u^{2}$，右端展开为 $u(2\pi^{2}u-\frac{(2\pi^{2})^{2}}{2}u^{2}+\cdots)$，首项同样是 $2\pi^{2}u^{2}$。这里的 $2\pi^{2}\kappa_{1}$ 与 §2.3 中称为 Weil–Petersson 类的 $2\pi^{2}\kappa_{1}$ 是同一个常数。]

\begin{theorem}[3.7（Mumford 公式）]
Hodge 类 $\Lambda(t)$ 是从平凡 CohFT 出发，依次经如下平移与旋转作用（按此顺序）而得到的 CohFT [Mum83]：
\begin{equation}
\begin{aligned}
  R(u)&=\exp\left(-\sum_{m\ge1}\frac{B_{m+1}}{m(m+1)}(tu)^{m}\right),\\
  T(u)&=u\big(1-R(u)\big),
\end{aligned}
  \tag{3.47}
\end{equation}
其中 $B_{m}$ 是第 $m$ 个 Bernoulli 数。在对稳定图求和进行重求和之后，我们推得
\begin{equation}
  \Lambda(t)=\exp\left(\sum_{m\ge1}\frac{B_{m+1}}{m(m+1)}t^{m}\left(\kappa_{m}-\sum_{i=1}^{n}\psi_{i}^{m}+\delta_{m}\right)\right),
  \tag{3.48}
\end{equation}
\end{theorem}
% [Tutor Note: (3.47) 的第二式确实定义了一个平移：由 $B_{2}=\frac16$ 得 $R(u)=1-\frac{t}{12}u+O(u^{2})$，故 $T(u)=u(1-R(u))=\frac{t}{12}u^{2}+O(u^{3})\in u^{2}\Q[[u]]$，与 (3.43) 的“在次数 $0$ 与 $1$ 处消失”吻合。另外，$m$ 为偶数（$m\ge2$）时 $m+1\ge3$ 为奇数，而 Bernoulli 数的奇数下标值全为零，故 $B_{m+1}=0$，于是实际只剩 $m=2k-1$ 的项，系数 $\frac{B_{2k}}{(2k-1)\,2k}$ 与常见的 Mumford 公式写法一致。]
% [Tutor Note: (3.48) 中的 $\delta_{m}$ 在原文 p.35 顶部（本译注本的下一块）才给出定义：$\delta_{m}=\frac12 j_{*}\sum_{k+\ell=m-1}\psi^{k}(\psi')^{\ell}$，其中 $j$ 是所有余维一边界层（单边稳定图）的并的嵌入，$\psi,\psi'$ 是结点处的两个 $\psi$ 类；同时那里的 $\Lambda(t)$ 即 Block F 中 (3.17) 的 $\Lambda(u)=\sum_{k=0}^{g}\lambda_{k}u^{k}$（Hodge 丛的 Chern 多项式）。核对 $m=1$ 项：$\delta_{1}=\frac12 j_{*}\mathbf{1}$ 是边界除子类的一半，而 (3.48) 的指数中 $m=1$ 项给出 $\lambda_{1}=\frac{1}{12}\left(\kappa_{1}-\sum_{i}\psi_{i}+\delta_{1}\right)$，即熟知的 Mumford 关系。]
